% Manuscript update: interaction between inherited and contemporary curvature
% Canonical manuscript source is maintained separately; this file records the
% interaction subsection added during the Milky Way persistence analysis.

\subsection{Interaction between inherited and contemporary curvature}
\label{sec:curvature-interaction}

The central physical hypothesis is stronger than simple superposition.  A
persistent geometric state generated by earlier baryonic evolution defines part
of the spacetime environment through which newly generated curvature propagates.
Consequently, the response associated with the contemporary baryonic source may
be realized in a geometry that is already displaced from the state that would
have been obtained had the earlier response relaxed completely.

We therefore allow the effective curvature to be represented schematically by a
covariant response functional
\begin{equation}
    G_{\mu\nu}^{\mathrm{eff}}
    =
    \mathcal{F}_{\mu\nu}
    \!\left[
        G_{\alpha\beta}^{(b)},
        M_{\alpha\beta}
    \right],
    \label{eq:nonlinear-effective-curvature}
\end{equation}
subject to the recovery condition
\begin{equation}
    \mathcal{F}_{\mu\nu}
    \!\left[
        G_{\alpha\beta}^{(b)},0
    \right]
    =
    G_{\mu\nu}^{(b)}.
    \label{eq:interaction-gr-limit}
\end{equation}
A weak-coupling expansion of \eqref{eq:nonlinear-effective-curvature} may then
be written as
\begin{equation}
    G_{\mu\nu}^{\mathrm{eff}}
    =
    G_{\mu\nu}^{(b)}
    +
    M_{\mu\nu}
    +
    \mathcal{C}_{\mu\nu}
    \!\left[
        G^{(b)},M
    \right]
    +
    \mathcal{O}(M^2,G^{(b)}M^2,\ldots),
    \label{eq:curvature-coupling-expansion}
\end{equation}
where \(\mathcal{C}_{\mu\nu}\) denotes the leading interaction between the
contemporary Einstein response and the inherited geometric state.  By
definition, the interaction term satisfies
\begin{equation}
    \mathcal{C}_{\mu\nu}[G^{(b)},0]=0,
    \qquad
    \mathcal{C}_{\mu\nu}[0,M]=0,
    \label{eq:coupling-vanishing}
\end{equation}
so that it represents genuine cross-coupling rather than a redefinition of
either sector alone.  Pure self-interactions of the memory sector, if required
by a nonlinear completion, are conceptually distinct and are absorbed into the
higher-order terms in \eqref{eq:curvature-coupling-expansion}.

This formulation distinguishes three effects: the ordinary instantaneous
Einstein response to the current stress--energy tensor, the surviving geometric
state deposited by prior baryonic configurations, and the modification produced
when contemporary curvature is generated within that inherited state.  The
observable gravitational residual therefore need not be identified with
\(M_{\mu\nu}\) alone.  In general it probes the combined hereditary contribution
\begin{equation}
    \Delta G_{\mu\nu}
    \equiv
    G_{\mu\nu}^{\mathrm{eff}}
    -
    G_{\mu\nu}^{(b)}
    =
    M_{\mu\nu}
    +
    \mathcal{C}_{\mu\nu}[G^{(b)},M]
    +\cdots .
    \label{eq:observable-geometric-residual}
\end{equation}
This distinction is essential observationally: two regions with comparable
present baryonic stress--energy may exhibit different effective gravitational
responses if their inherited geometric states differ, while two regions with
similar inherited states may respond differently if their present baryonic
configurations differ.  These matched-history and matched-present-state
comparisons provide a direct empirical test of whether the persistence sector
acts additively or through a nontrivial interaction with contemporary curvature.
